\chapter{Bagian Umum}

\section{Tabel kontribusi}

Tabel ini menunjukkan siapa mengerjakan apa. Pembagian kerja mengikuti tiga tahap: fondasi dan Problem~1, lalu
Problem~2 dan~3, lalu Problem~4 dan~5 beserta integrasi.

\begin{table}[!htbp]
\centering
\small
\begin{tabularx}{\textwidth}{@{}L{2.9cm}YYY@{}}
\toprule
\textbf{Anggota} & \textbf{Perancangan} & \textbf{Implementasi} & \textbf{Penulisan laporan}\\
\midrule
Orang~1 \newline Andi Farhan Hidayat &
 Kontrak antar-komponen, arsitektur Aggregator, kebijakan skema &
 Fondasi (\code{pkg/obs}, Compose, Makefile, CI), mock BMKG dan PVMBG, Aggregator, Problem~1, ketahanan ingest, adapter PostgreSQL dan Kafka, outbox dan relay, consumer &
 Bagian umum, Problem~1, deskripsi mock\\
Orang~2 \newline Muhammad Rizain Firdaus &
 Autentikasi JWT tiga peran, kebijakan field Ringkasan/Mentah, pembatas laju, penanda \code{stale} &
 \code{auth-service}, \code{client-api}, skrip demo dan load test Problem~2 dan~3 &
 Problem~2, Problem~3\\
Orang~3 \newline Muhammad Kinan Arkansyaddad &
 Perbandingan pilihan penyimpanan (PostgreSQL dengan JSONB) dan pilihan penyebaran event (Kafka dengan outbox), semantik pengiriman, dan batasan yang disadari &
 Skrip demo Problem~4 dan~5 (membangun ulang satu service, field baru tanpa migrasi, consumer baru) dan pengujian pemulihan consumer &
 Problem~4, Problem~5\\
\bottomrule
\end{tabularx}
\caption{Kontribusi anggota.}
\end{table}

\section{Deskripsi umum sistem}

Di cerita soal, Astolfo menerima sinyal bencana dari banyak lembaga yang tidak saling bicara. Sistem ini, ARGUS
(\emph{Aggregated Real-time Geohazard Unified System}), adalah \emph{gateway} milik BNPB yang menyambungkannya. Ia melakukan lima hal:

\begin{enumerate}
 \item \textbf{Mengambil} data dari dua lembaga secara berkala (\emph{polling}): BMKG untuk gempa dan tsunami, PVMBG untuk gunung api.
 \item \textbf{Menyeragamkan} dua format yang berbeda menjadi satu format milik BNPB, yaitu \emph{HazardEvent}.
 \item \textbf{Menyimpan} hasilnya di satu penyimpanan, \emph{Canonical Store}, yang hanya dikelola oleh satu service.
 \item \textbf{Menyebarkan} setiap data baru ke pihak yang berlangganan lewat \emph{message broker}.
 \item \textbf{Menyajikan} data ke tiga jenis klien (Media, Tim Lapangan, BNPB Pusat) dengan hak akses yang berbeda.
\end{enumerate}

\textbf{Lingkup Milestone 1} dibatasi oleh soal: dua lembaga sumber (BMKG dan PVMBG), dua jenis bencana (seismik dan
vulkanik), BNPB sebagai simpul koordinasi, dan semua klien hanya membaca (\emph{read-only}). BMKG dan PVMBG adalah
\emph{mock} buatan kelompok yang mengikuti kontrak di soal tanpa mengubahnya.

\subsection*{Status penyelesaian}

Tabel berikut merangkum apa yang sudah terbukti lewat demo. Bukti P1 dan sisi ingest P2 dibuat pada 7~Oktober~2026;
bukti sisi layanan P2, P3, P4, dan P5 pada 8~Oktober~2026. Demo P2 dan P3 dijalankan sekali setelah \cmd{make reset}.
Demo P4 dan P5 dijalankan setelah reset dan diulang sekali tanpa reset; keduanya lulus kedua kali.

\begin{table}[!htbp]
\centering
\small
\begin{tabularx}{\textwidth}{@{}L{3.0cm}Y L{2.0cm}@{}}
\toprule
\textbf{Problem} & \textbf{Yang sudah terbukti $\;/\;$ yang masih kurang} & \textbf{Status}\\
\midrule
P1 Interoperabilitas data & Pemetaan benar, skema PVMBG berubah saat berjalan tanpa restart, field baru masuk \code{attributes}, perubahan yang tidak didukung dikarantina. & \lulus\\
P2 Konkurensi dan ketersediaan & p95 BMKG-only di bawah 300~ms saat PVMBG 3~detik, load test 50~koneksi selama 60~detik, outage PVMBG ditangani, pulih tanpa restart. & \lulus\\
P3 Autentikasi dan kepercayaan & Kredensial antar-lembaga saling menolak, Media hanya Ringkasan, token Lapangan kedaluwarsa lalu di-\emph{refresh}, tidak ada kredensial ter-commit. & \lulus\\
P4 Microservice dan storage & Satu container (\code{pvmbg-mock}) dihentikan, dibangun ulang, dan dijalankan lagi tanpa menyentuh yang lain; field baru masuk JSONB tanpa migrasi dan record lama tetap terbaca; Canonical Store hanya terjangkau oleh Aggregator. & \lulus\\
P5 Pub/sub & Event lewat Kafka, dua consumer independen, satu consumer dimatikan lalu menyusul event yang tepat, dan consumer ketiga ditambahkan saat sistem berjalan tanpa mengubah producer. & \lulus\\
\bottomrule
\end{tabularx}
\caption{Status per problem pada saat laporan ini ditulis.}
\end{table}

\section{Diagram arsitektur}

Gambar~\ref{fig:arsitektur} memperlihatkan seluruh service, protokol, kepemilikan penyimpanan, dan batas jaringan.

% Gambar 1.1 -- arsitektur sistem. Dipanggil dari umum.tex.
\begin{figure}[!htbp]
\centering
\resizebox{\textwidth}{!}{%
\begin{tikzpicture}[
  font=\sffamily\footnotesize,
  svc/.style={draw=accent, fill=white, rounded corners=3pt, thick, align=center, minimum height=1.05cm, inner sep=4pt},
  mock/.style={svc, draw=gridgray, fill=gray!8},
  store/.style={draw=okgreen, fill=okbg, thick, cylinder, shape border rotate=90, aspect=0.25, align=center, minimum height=1.4cm, minimum width=2.6cm},
  zone/.style={draw=gridgray, dashed, rounded corners=6pt, inner sep=8pt},
  flow/.style={-{Stealth[length=2.2mm]}, thick, color=ink},
  dotflow/.style={-{Stealth[length=2.2mm]}, thick, dotted, color=ink},
  lab/.style={font=\sffamily\scriptsize, color=ink, fill=white, inner sep=1.5pt, align=center}
]
% ---- instansi sumber
\node[mock] (bmkg) at (0,6.4) {\textbf{bmkg-mock}\\:8081 \;\texttt{X-BMKG-Key}};
\node[mock] (pvmbg) at (0,3.6) {\textbf{pvmbg-mock}\\:8082 \;\texttt{Bearer pvmbg.v1...}};
% ---- aggregator
\node[svc, fill=accentlight, minimum width=3.9cm, minimum height=3.1cm] (agg) at (6.4,5.0)
  {\textbf{aggregator}\\[2pt]polling per sumber\\pemetaan ke HazardEvent\\outbox + relay\\API internal :8083};
% ---- sisi klien
\node[svc] (client) at (13.4,8.6) {\textbf{Klien}\\Media $\cdot$ Lapangan $\cdot$ Pusat};
\node[svc] (auth) at (9.2,8.6) {\textbf{auth-service}\\:8084};
\node[svc] (capi) at (13.4,5.6) {\textbf{client-api}\\:8080};
% ---- penyimpanan dan broker
\node[store] (pg) at (6.4,0.4) {\textbf{PostgreSQL}\\Canonical Store};
\node[svc] (kafka) at (10.6,1.7) {\textbf{Kafka}\\topic \texttt{hazard.events}};
\node[svc] (c1) at (14.0,2.9) {\textbf{consumer-dashboard}\\:8091};
\node[svc] (c2) at (14.0,0.5) {\textbf{consumer-notifier}\\:8092};
% ---- zona jaringan
\begin{scope}[on background layer]
 \node[zone, fit=(pg), label={[font=\sffamily\scriptsize,color=gridgray]below:canonical-net (internal)}] {};
 \node[zone, fit=(kafka)(c1)(c2), label={[font=\sffamily\scriptsize,color=gridgray]below:broker-net (internal)}] {};
\end{scope}
% ---- panah
\draw[flow] (agg.west |- bmkg.east) -- (bmkg.east) node[lab, pos=0.45] {poll HTTP};
\draw[flow] (agg.west |- pvmbg.east) -- (pvmbg.east) node[lab, pos=0.45] {poll HTTP};
\draw[flow] (agg.south) -- (pg.north) node[lab, pos=0.5] {baca / tulis};
\draw[flow] (agg.south east) -- (kafka.west) node[lab, pos=0.45] {outbox $\rightarrow$ publish};
\draw[flow] (kafka.east) -- (c2.west);
\draw[flow] (kafka.east) -- (c1.west);
\draw[flow] (client) -- (capi) node[lab, pos=0.5] {Bearer JWT};
\draw[flow] (client) -- (auth) node[lab, pos=0.5] {token / refresh};
\draw[dotflow] (auth.south east) -- (capi.north west) node[lab, pos=0.5] {kunci JWT\\sama};
\draw[flow] (capi.west) -- (capi.west -| agg.east) node[lab, pos=0.5] {\texttt{X-Service-Key}};
\end{tikzpicture}}
\caption{Arsitektur sistem. Garis tegas: panggilan jaringan; garis titik: rahasia penandatangan token yang dipakai bersama. Kotak putus-putus: jaringan Docker yang tidak bisa dijangkau dari luar (\texttt{internal}).}
\label{fig:arsitektur}
\end{figure}


\subsection*{Penjelasan}

\textbf{Alur masuk.} \code{aggregator} memanggil \code{bmkg-mock} dan \code{pvmbg-mock} lewat HTTP. Kedua mock pasif: mereka
hanya menjawab, tidak mengirim data sendiri. Setiap sumber punya putaran polling sendiri, sehingga PVMBG yang lambat
tidak menahan BMKG (Problem~2). Hasil polling dipetakan ke \emph{HazardEvent} dan disimpan.

\textbf{Penyimpanan.} Canonical Store adalah PostgreSQL yang hanya berada di jaringan \code{canonical-net}. Jaringan itu
hanya berisi \code{aggregator} dan database-nya, jadi service lain secara fisik tidak bisa menyambung ke database
(Problem~4). Service lain memperoleh data lewat API internal Aggregator.

\textbf{Alur keluar ke klien.} Klien meminta token ke \code{auth-service}, lalu memanggil \code{client-api} dengan token
itu. \code{client-api} memeriksa token, menentukan field apa yang boleh dilihat peran tersebut, dan meneruskan
permintaan ke API internal Aggregator memakai \code{X-Service-Key}. \code{client-api} tidak pernah memanggil BMKG
atau PVMBG (Problem~2 dan~3).

\textbf{Alur event.} Saat Aggregator menyimpan perubahan, ia menulis juga satu baris \emph{outbox} dalam transaksi
database yang sama. Proses \emph{relay} lalu mengirim baris itu ke topic Kafka \code{hazard.events}. Consumer
berlangganan ke topic tersebut tanpa diketahui Aggregator (Problem~5). Setiap alur membawa \emph{correlation ID} yang
sama dari siklus polling sampai ke consumer.

\subsection*{Daftar container}

Sistem terdiri dari sepuluh container, satu di antaranya hanya tugas sekali jalan. Tabel berikut merangkum port, jaringan, dan penyimpanan masing-masing.

\begin{table}[!htbp]
\centering
\small
\begin{tabularx}{\textwidth}{@{}L{3.7cm}L{2.5cm}L{2.5cm}Y@{}}
\toprule
\textbf{Container} & \textbf{Port host} & \textbf{Jaringan} & \textbf{Penyimpanan yang dimiliki}\\
\midrule
\code{bmkg-mock} & 8081 & app-net & tidak ada (data simulasi di memori)\\
\code{pvmbg-mock} & 8082 & app-net & tidak ada (data simulasi di memori)\\
\code{aggregator} & 127.0.0.1:8083 & app, canonical, broker & Canonical Store (tabel \code{hazards}, \code{outbox}, \code{poll_cursor})\\
\code{postgres-canonical} & tidak dibuka & canonical-net & volume \code{canonical-pgdata}\\
\code{kafka} & tidak dibuka & broker-net & volume \code{kafka-data}\\
\code{kafka-init} & -- & broker-net & tugas sekali jalan: membuat topic\\
\code{auth-service} & 8084 & app-net, auth-net & SQLite untuk hash refresh token (volume \code{auth-data})\\
\code{client-api} & 8080 & app-net & tidak ada (tanpa \emph{state})\\
\code{consumer-dashboard} & 8091 & broker, app & memori (himpunan \code{event_id})\\
\code{consumer-notifier} & 8092 & broker, app & memori (himpunan \code{event_id})\\
\bottomrule
\end{tabularx}
\caption{Container, jaringan, dan kepemilikan penyimpanan (dari \texttt{docker-compose.yml}).}\label{tab:container}
\end{table}

\section{Teknologi terpilih}

Tabel berikut menjawab pertanyaan: teknologi apa yang dipakai untuk tiap komponen, dan kebutuhan sistem mana yang membuatnya dipilih.

\begin{table}[!htbp]
\centering
\small
\begin{tabularx}{\textwidth}{@{}L{3.3cm}L{3.8cm}Y@{}}
\toprule
\textbf{Komponen} & \textbf{Teknologi} & \textbf{Alasan (dikaitkan dengan kebutuhan)}\\
\midrule
Semua service & Go 1.25 & Goroutine ringan cocok untuk satu putaran polling per sumber. Image kecil dan build cepat, sehingga membangun ulang satu service tidak lama (Problem~4). Satu bahasa menekan biaya integrasi tiga orang.\\
HTTP dan log & \code{net/http}, \code{log/slog} & Tanpa dependensi tambahan. Log JSON membawa \code{cid} lewat \emph{context}, sesuai syarat log terstruktur.\\
Canonical Store & PostgreSQL 16 dengan kolom JSONB & Kolom tetap untuk field kanonik, \code{attributes} JSONB untuk field yang bisa berubah: field baru tidak butuh migrasi (Problem~1 dan~4). Transaksi memungkinkan pola outbox (Problem~5).\\
Message broker & Apache Kafka (mode KRaft, satu broker) & Log tahan lama dengan \emph{consumer group} dan \emph{offset}: consumer yang mati dapat menyusul setelah hidup lagi (Problem~5).\\
Klien Kafka & \code{franz-go} & Go murni, tanpa cgo, sehingga Dockerfile sederhana.\\
Autentikasi & JWT HS256 dan refresh token buram & Token pendek yang diperiksa lokal, jadi jalur baca tidak bergantung pada Auth Service (Problem~2 dan~3).\\
Penyimpanan token & SQLite (\code{modernc.org/sqlite}) & Hanya hash refresh token yang disimpan; cukup untuk satu instance dan tanpa cgo.\\
Pembatas laju & \code{golang.org/x/time/rate} & Token bucket: menolak kelebihan beban seketika dengan 429, bukan mengantre (Problem~2).\\
Orkestrasi & Docker Compose, empat jaringan & Satu perintah menyalakan semua. Jaringan \emph{internal} menjadi bukti isolasi penyimpanan.\\
Load test & k6 & Melaporkan throughput, p50/p95/p99, dan laju galat; mendukung ambang lulus.\\
Kualitas & GitHub Actions, \code{gofmt}, \cmd{go vet}, \code{-race}, gitleaks & Tes otomatis dan pencegahan kredensial ter-commit.\\
\bottomrule
\end{tabularx}
\caption{Teknologi per komponen dan alasan pemilihannya.}
\end{table}

\section{Asumsi}

\begin{enumerate}
 \item Respons mock berupa JSON array. Parameter \code{since} bersifat \emph{lebih besar dari} (bukan sama dengan) terhadap \code{occurred_at}, \code{issued_at}, atau \code{reported_at}. Field \code{issued_at} pada TsunamiWarning adalah tambahan yang diizinkan soal.
 \item Warning tsunami dapat terbit \emph{setelah} gempanya. Mock sengaja memberi jeda 2--8 detik. Satu gempa dapat punya beberapa warning dan yang tertinggi menentukan severity.
 \item Hanya ada satu instance tiap service pada M1. Konkurensi optimistis (\code{expectedVersion}) sudah disiapkan untuk M2.
 \item Koordinat dan nama gunung api berasal dari tabel referensi statis milik BNPB (\code{config/volcanoes.json}). Koordinatnya pendekatan untuk simulasi, dan gunung ``Prikitiw'' fiktif mengikuti cerita soal.
 \item Seluruh waktu memakai UTC dan dipotong ke ketelitian mikrodetik.
 \item Karantina record yang tidak dapat dipetakan bersifat diagnostik dan disimpan di memori.
 \item Endpoint \code{/admin/*} pada mock (termasuk \code{extra_fields}, \code{override}, \code{/admin/generate}, \code{/admin/delay}) adalah tambahan untuk demo. Kontrak yang ditetapkan soal tidak diubah.
 \item Semua klien hanya membaca. Jalur pelaporan dari lapangan ke BNPB berada di luar M1.
 \item Jaringan Docker dianggap tepercaya; tidak ada TLS di dalamnya. Pada produksi, TLS dipasang di ingress.
 \item Kredensial pada \code{.env.example} hanyalah contoh. Nilai sebenarnya diisi di \code{.env} yang tidak pernah ter-commit.
\end{enumerate}

\section{Petunjuk menjalankan}

Ringkasan di sini mengacu ke \code{README.md} di repositori, yang memuat langkah lengkap, daftar variabel, dan cara
memicu setiap skenario.

\begin{lstlisting}
cp .env.example .env            # nilai contoh; .env tidak pernah di-commit
docker compose up --build       # satu perintah untuk seluruh sistem (atau: make up)
make demo-p1                    # lalu demo-p2, demo-p3, demo-isolation, demo-p4, demo-p5, dst.
make reset                      # kembali ke kondisi bersih (hapus container dan volume)
\end{lstlisting}

Satu-satunya prasyarat adalah Docker dengan Compose v2. Port bawaan: BMKG 8081, PVMBG 8082, Client-Facing API 8080,
Auth Service 8084, Aggregator 8083 (hanya \emph{loopback}), consumer 8091 dan 8092. Semua dapat diubah lewat
\code{.env}. Setiap demo memeriksa dirinya sendiri (kode keluar tidak nol bila gagal) dan menyimpan transkripnya ke
\code{docs/evidence/}.

\section{Istilah penting}

Beberapa istilah muncul berulang di laporan ini. Artinya dirangkum di sini agar mudah dicari.

\begin{table}[!htbp]
\centering
\small
\begin{tabularx}{\textwidth}{@{}L{3.4cm}Y@{}}
\toprule
\textbf{Istilah} & \textbf{Artinya dalam laporan ini}\\
\midrule
HazardEvent & Format tunggal milik BNPB untuk satu kejadian bahaya, apa pun sumbernya.\\
Canonical Store & Tempat HazardEvent disimpan. Hanya Aggregator yang boleh membukanya langsung.\\
Ringkasan / Mentah & Dua klasifikasi field. Ringkasan boleh dilihat Media; Mentah (koordinat presisi, \code{attributes}, \code{source_ref_id}) tidak.\\
Polling & BNPB yang bertanya ke sumber secara berkala. Sumber tidak mengirim sendiri.\\
Outbox & Tabel di database yang menampung event sampai terkirim ke broker, supaya tidak ada event hilang bila broker sedang mati.\\
At-least-once & Setiap event pasti sampai, tetapi bisa sampai lebih dari sekali.\\
Idempoten & Memproses hal yang sama dua kali memberi hasil yang sama dengan sekali.\\
Circuit breaker & Pemutus: setelah sumber gagal berkali-kali, panggilan dihentikan sementara agar tidak memboroskan sumber daya.\\
Token bucket & Pembatas laju: tiap permintaan mengambil satu token; token habis berarti ditolak dengan 429.\\
Stale & Data masih ditampilkan tetapi diberi tanda bahwa sumbernya sedang tidak memperbarui.\\
Correlation ID (\code{cid}) & Kode yang sama pada semua log satu alur permintaan, dari service pertama sampai terakhir.\\
\bottomrule
\end{tabularx}
\caption{Daftar istilah.}
\end{table}
